%% kapa: exam=1 round=19 year=2008 period=1 form=A subject=civil_law subject_ko=민법 seq=30 number=30
%% source: Q-Net 감정평가사 기출문제 (cst00309) artlSeq=1001962 file=extracted/1008/1001962/2008년 제19회 감정평가사 1차시험 기출문제/1교시.hwp
%% conversion: kordoc 4.24.51 -> tools/gen1.py (auto); review=proofread-sonnet; flags=none
\begin{question}{30}
지상권에 관한 설명으로 틀린 것은? (다툼이 있는 경우에는 판례에 의함)
\begin{choices}
  \item 지상권이 존속기간의 만료로 소멸한 경우에 건물 기타 공작물이나 수목이 현존하고 있으면 지상권자는 계약의 갱신을 청구할 수 있으며, 지상권설정자는 이를 거절할 수 없다.
  \item 토지소유자는 법정지상권이 성립한 후에 증축된 건물부분에 대하여 그 철거를 청구할 수 없다.
  \item 구분지상권에 의하여 소유할 수 있는 목적물에는 일반지상권과 달리 수목은 포함되지 않는다.
  \item 당사자가 계약으로 지상권의 존속기간을 약정하지 않아도 지상권은 유효하게 성립할 수 있다.
  \item 지상권자가 2년 이상의 지료지급을 지체한 경우, 지상권설정자는 지료지급이 지체되었음을 이유로 지상권의 소멸을 청구할 수 있다.
\end{choices}
\end{question}
\answer{1}
